\chapter{文件系统}
\label{CH:FS}
%         这是进程还是内核线程？
% 
%         日志文本（以及部分缓冲区文本）假设读者对 inode 和目录
%         的工作方式有相当了解，但它们将在后面介绍。
% 
%         必须决定使用 processor 还是 CPU。
% 
%         确保说的是 buffer，不是 block
% 
%         Mount

文件系统的目的是组织和存储数据。文件系统通常支持用户和应用程序之间共享数据，并支持持久性，以便重启后数据仍然可用。

xv6 文件系统提供了类似 Unix 的文件、目录和路径名（见第~\ref{CH:UNIX}~章），并将其数据存储在 virtio 磁盘上以实现持久性。该文件系统应对几个挑战：
\begin{itemize}
  
\item 文件系统需要磁盘上的数据结构来表示命名的目录和文件树，记录保存每个文件内容的块的标识，并记录磁盘的哪些区域是空闲的。
\item 文件系统必须支持崩溃恢复。也就是说，如果发生崩溃（例如断电），文件系统在重启后必须仍然正确工作。风险在于崩溃可能中断一系列更新，并导致磁盘上的数据结构不一致（例如，一个块既在文件中使用又被标记为空闲）。
\item 不同的进程可能同时操作文件系统，因此文件系统代码必须进行协调以维护不变量。
\item 访问磁盘比访问内存慢几个数量级，因此文件系统必须维护一个常用块的内存缓存。

\end{itemize}

本章其余部分解释 xv6 如何应对这些挑战。
%% 
%%  -------------------------------------------
%% 
\section{概述}

xv6 文件系统实现分为七层，如图~\ref{fig:fslayer}~所示。磁盘层在 virtio 硬盘上读写块。缓冲区缓存层缓存磁盘块并同步对它们的访问，确保每次只有一个内核进程可以修改任何特定块中存储的数据。日志层允许上层将对多个块的更新包装在一个事务中，并确保在崩溃时这些块被原子性地更新（即要么全部更新，要么都不更新）。inode 层提供单个文件，每个文件表示为一个具有唯一 i-number 和一些保存文件数据块的 inode。目录层将每个目录实现为一种特殊类型的 inode，其内容是一系列目录项，每个目录项包含一个文件名和 i-number。路径名层提供分层路径名，如 \lstinline{/usr/rtm/xv6/fs.c}，并通过递归查找解析它们。文件描述符层使用文件系统接口抽象许多 Unix 资源（例如管道、设备、文件等），简化了应用程序员的生活。

\begin{figure}[t]
\center
\includegraphics[scale=0.5]{fig/fslayer.pdf}
\caption{xv6 文件系统的层次。}
\label{fig:fslayer}
\end{figure}

传统上，磁盘硬件将磁盘上的数据显示为编号的 512 字节块序列（也称为扇区）：扇区 0 是前 512 字节，扇区 1 是下一个，依此类推。操作系统为其文件系统使用的块大小可能与磁盘使用的扇区大小不同，但通常块大小是扇区大小的倍数。Xv6 将它读入内存的块的副本保存在 \lstinline{struct buf} 类型的对象中。此结构中存储的数据有时与磁盘不同步：它可能尚未从磁盘读入（磁盘正在处理但尚未返回扇区内容），或者可能已被软件更新但尚未写入磁盘。

文件系统必须有一个计划来存储 inode 和内容块在磁盘上的位置。为此，xv6 将磁盘划分为几个部分，如图~\ref{fig:fslayout}~所示。文件系统不使用块 0（它包含引导扇区）。块 1 称为超级块；它包含有关文件系统的元数据（以块为单位的文件系统大小、数据块数量、inode 数量以及日志中的块数量）。从块 2 开始的块保存日志。日志之后是 inode，每个块有多个 inode。在这些之后是位图块，用于跟踪哪些数据块正在使用。其余的是数据块；每个数据块在位图块中要么被标记为空闲，要么保存文件或目录的内容。超级块由一个单独的程序 \indexcode{mkfs} 填充，该程序构建初始文件系统。

本章的其余部分将从缓冲区缓存开始讨论每一层。请注意较低层中精心选择的抽象如何简化较高层设计的场景。
%% 
%%  -------------------------------------------
%% 
\section{缓冲区缓存层}
\label{s:bcache}

缓冲区缓存有两个任务：(1) 同步对磁盘块的访问，确保内存中只有一个块的副本，并且每次只有一个内核线程使用该副本；(2) 缓存常用块，这样就不必从慢速磁盘重新读取它们。代码在 \lstinline{bio.c} 中。

缓冲区缓存导出的主要接口是 \indexcode{bread} 和 \indexcode{bwrite}；前者获取一个包含块副本的 \indextext{buf}，可以在内存中读取或修改，后者将修改后的缓冲区写入磁盘上的相应块。内核线程在使用完缓冲区后必须调用 \indexcode{brelse} 释放它。缓冲区缓存使用每个缓冲区的睡眠锁来确保每次只有一个线程使用每个缓冲区（以及每个磁盘块）；\lstinline{bread} 返回一个已锁定的缓冲区，\lstinline{brelse} 释放锁。

让我们回到缓冲区缓存。缓冲区缓存有固定数量的缓冲区来保存磁盘块，这意味着如果文件系统请求一个尚未在缓存中的块，缓冲区缓存必须回收当前保存其他某个块的缓冲区。缓冲区缓存为新的块回收最近最少使用的缓冲区。假设是最近最少使用的缓冲区是未来最不可能再次使用的。

\begin{figure}[t]
\center
\includegraphics[scale=0.5]{fig/fslayout.pdf}
\caption{xv6 文件系统的结构。}
\label{fig:fslayout}
\end{figure}
%% 
%%  -------------------------------------------
%% 
\section{代码：缓冲区缓存}

缓冲区缓存是一个双向链表。函数 \indexcode{binit} 由 \indexcode{main} 调用，用静态数组 \lstinline{buf} 中的 \indexcode{NBUF} 个缓冲区初始化链表。对缓冲区缓存的所有其他访问都通过 \indexcode{bcache.head} 引用链表，而不是 \lstinline{buf} 数组。

一个缓冲区有两个相关的状态字段。字段 \indexcode{valid} 指示缓冲区包含块的副本。字段 \indexcode{disk} 指示缓冲区内容已交给磁盘，磁盘可能更改缓冲区（例如，将数据从磁盘写入 \lstinline{data}）。

\lstinline{bread} 调用 \indexcode{bget} 来获取给定扇区的缓冲区。如果缓冲区需要从磁盘读取，\lstinline{bread} 在返回缓冲区之前调用 \indexcode{virtio_disk_rw} 来执行读取。

\lstinline{bget} 扫描缓冲区列表，查找具有给定设备和扇区号的缓冲区。如果存在这样的缓冲区，\indexcode{bget} 获取该缓冲区的睡眠锁。\lstinline{bget} 然后返回锁定的缓冲区。

如果没有给定扇区的缓存缓冲区，\indexcode{bget} 必须创建一个，可能重用持有不同扇区的缓冲区。它第二次扫描缓冲区列表，寻找一个未使用的缓冲区（\lstinline{b->refcnt = 0}）；任何这样的缓冲区都可以使用。\lstinline{bget} 编辑缓冲区元数据以记录新的设备和扇区号，并获取其睡眠锁。注意赋值 \lstinline{b->valid = 0} 确保 \lstinline{bread} 将从磁盘读取块数据，而不是错误地使用缓冲区以前的内容。

每个磁盘扇区最多有一个缓存的缓冲区很重要，以确保读者看到写入，并且因为文件系统使用缓冲区上的锁进行同步。\lstinline{bget} 通过连续持有 \lstinline{bcache.lock} 来保证这个不变性，从第一个循环检查块是否被缓存，一直到第二个循环声明该块现在已被缓存（通过设置 \lstinline{dev}、\lstinline{blockno} 和 \lstinline{refcnt}）。这使得检查块是否存在以及（如果不存在）指定一个缓冲区来保存该块成为原子操作。

在 \lstinline{bcache.lock} 临界区之外获取缓冲区的睡眠锁是安全的，因为非零的 \lstinline{b->refcnt} 阻止该缓冲区被重用于不同的磁盘块。睡眠锁保护对缓冲区内容的读写，而 \lstinline{bcache.lock} 保护有关哪些块被缓存的信息。

如果所有缓冲区都忙，则意味着太多进程同时执行文件系统调用；\lstinline{bget} 会 panic。更优雅的响应可能是休眠直到有缓冲区可用，但那样可能会产生死锁的可能性。

一旦 \indexcode{bread} 读取了磁盘（如果需要）并将缓冲区返回给其调用者，调用者就独占了该缓冲区，可以读取或写入数据字节。如果调用者修改了缓冲区，它必须在释放缓冲区之前调用 \indexcode{bwrite} 将更改的数据写入磁盘。\lstinline{bwrite} 调用 \indexcode{virtio_disk_rw} 与磁盘硬件通信。

当调用者使用完缓冲区后，它必须调用 \indexcode{brelse} 来释放它。（名称 \lstinline{brelse} 是 b-release 的缩写，虽然神秘但值得学习：它起源于 Unix，并在 BSD、Linux 和 Solaris 中使用。）\lstinline{brelse} 释放睡眠锁，并将缓冲区移到链表的前面。移动缓冲区使得链表按缓冲区最近使用（即释放）的顺序排列：链表中的第一个缓冲区是最近使用的，最后一个是最近最少使用的。\lstinline{bget} 中的两个循环利用了这一点：在最坏情况下，查找现有缓冲区的扫描必须处理整个链表，但从最近使用的缓冲区开始检查（从 \lstinline{bcache.head} 开始并跟随 \lstinline{next} 指针）可以在有良好局部性引用时减少扫描时间。选择要重用缓冲区的扫描通过向后扫描（跟随 \lstinline{prev} 指针）来选择最近最少使用的缓冲区。
%% 
%%  -------------------------------------------
%% 
\section{日志层}

文件系统设计中最有趣的问题之一是崩溃恢复。问题的出现是因为许多文件系统操作涉及对磁盘的多次写入，而在一部分写入之后发生的崩溃可能使磁盘上的文件系统处于不一致的状态。例如，假设在文件截断期间（将文件长度设置为零并释放其内容块）发生崩溃。根据磁盘写入的顺序，崩溃可能要么留下一个引用被标记为空闲的内容块的 inode，要么留下一个已分配但未被引用的内容块。

后者相对无害，但引用已释放块的 inode 很可能在重启后引起严重问题。重启后，内核可能会将该块分配给另一个文件，现在我们有两个不同的文件无意中指向同一个块。如果 xv6 支持多用户，这种情况可能是一个安全问题，因为旧文件的所有者将能够读写新文件中的块，而新文件属于另一个用户。

Xv6 用一种简单的日志形式解决了文件系统操作期间崩溃的问题。一个 xv6 系统调用不会直接写入磁盘上的文件系统数据结构。相反，它将它希望执行的所有磁盘写入的描述放在磁盘上的一个日志中。一旦系统调用记录了其所有写入，它就会在磁盘上写入一个特殊的提交记录，指示日志包含一个完整的操作。此时，系统调用将写入复制到磁盘上的文件系统数据结构。在这些写入完成后，系统调用会擦除磁盘上的日志。

如果系统崩溃并重新启动，文件系统代码在运行任何进程之前会按如下方式从崩溃中恢复。如果日志被标记为包含完整的操作，则恢复代码将写入复制到磁盘上文件系统中它们所属的位置。如果日志未被标记为包含完整的操作，则恢复代码忽略该日志。恢复代码最后擦除日志。

为什么 xv6 的日志解决了文件系统操作期间的崩溃问题？如果在操作提交之前发生崩溃，则磁盘上的日志不会被标记为完成，恢复代码将忽略它，磁盘的状态就像操作从未开始一样。如果在操作提交之后发生崩溃，则恢复将重放该操作的所有写入，即使操作已经开始将它们写入磁盘上的数据结构，也可能会重复它们。无论哪种情况，日志都使操作在崩溃方面具有原子性：恢复后，操作的所有写入要么都出现在磁盘上，要么都不出现。
%% 
%% 
%% 
\section{日志设计}

日志位于超级块中指定的已知固定位置。它由一个头块和一系列更新的块副本组成。头块包含一个扇区号数组，每个日志块对应一个，以及日志块的数量。磁盘上头块中的计数要么为零，表示日志中没有事务，要么为非零，表示日志包含一个完整的已提交事务，其中包含指定数量的日志块。Xv6 在事务提交时写入头块，但在此之前不写入，并在将日志块复制到文件系统后将计数设置为零。因此，在事务中途发生的崩溃将导致日志头块中的计数为零；提交后发生的崩溃将导致非零计数。

每个系统调用的代码指示必须原子性应对崩溃的写入序列的开始和结束。为了允许不同进程并发执行文件系统操作，日志系统可以将多个系统调用的写入累积到一个事务中。因此，单个提交可能涉及多个完整系统调用的写入。为了避免将系统调用拆分到多个事务中，日志系统仅在没有文件系统系统调用正在进行时才提交。

将多个事务一起提交的想法称为组提交。组提交减少了磁盘操作的数量，因为它将提交的固定成本分摊到多个操作上。组提交还可以同时向磁盘系统提供更多并发写入，或许允许磁盘在单次磁盘旋转期间将它们全部写入。Xv6 的 virtio 驱动程序不支持这种批处理，但 xv6 的文件系统设计允许它。

Xv6 在磁盘上预留固定大小的空间来保存日志。事务中系统调用写入的块总数必须适合该空间。这有两个后果。不允许任何单个系统调用写入的不同的块数超过日志中的空间。这对于大多数系统调用来说不是问题，但其中两个调用可能写入许多块：\indexcode{write} 和 \indexcode{unlink}。大文件写入可能写入许多数据块和许多位图块以及一个 inode 块；取消链接一个大文件可能写入许多位图块和一个 inode。Xv6 的 write 系统调用将大的写入分解为多个适合日志的较小写入，而 \lstinline{unlink} 不会引起问题，因为实际上 xv6 文件系统只使用一个位图块。有限日志空间的另一个后果是，日志系统不能允许系统调用启动，除非确定该调用的写入将适合日志中剩余的空间。
%% 
%% 
%% 
\section{代码：日志记录}

系统调用中日志的典型用法如下所示：
\begin{lstlisting}[]
  begin_op();
  ...
  bp = bread(...);
  bp->data[...] = ...;
  log_write(bp);
  ...
  end_op();
\end{lstlisting}

\indexcode{begin_op} 等待直到日志系统当前不在提交，并且有足够的未预留日志空间来容纳此次调用的写入。\lstinline{log.outstanding} 统计已预留日志空间的系统调用数量；总预留空间是 \lstinline{log.outstanding} 乘以 \lstinline{MAXOPBLOCKS}。增加 \lstinline{log.outstanding} 既预留了空间，又阻止了在此系统调用期间发生提交。代码保守地假设每个系统调用可能写入最多 \lstinline{MAXOPBLOCKS} 个不同的块。

\indexcode{log_write} 充当 \indexcode{bwrite} 的代理。它记录块的扇区号在内存中，在磁盘上的日志中为其预留一个槽位，并将缓冲区固定在块缓存中，以防止块缓存将其驱逐。该块必须保留在缓存中直到提交：在此之前，缓存副本是修改的唯一记录；在提交之后才能将其写入磁盘上的位置；并且同一事务中的其他读取必须看到修改。\lstinline{log_write} 会注意到在单个事务中多次写入一个块的情况，并为该块分配日志中相同的槽位。这种优化通常称为吸收。例如，在单个事务中多次写入包含多个文件 inode 的磁盘块是很常见的。通过将多次磁盘写入吸收为一次，文件系统可以节省日志空间并获得更好的性能，因为磁盘块只需写入磁盘一次。

\indexcode{end_op} 首先减少未完成的系统调用计数。如果计数现在为零，它通过调用 \lstinline{commit()} 提交当前事务。这个过程有四个阶段。\lstinline{write_log()} 将事务中修改的每个块从缓冲区缓存复制到磁盘上日志中其对应的槽位。\lstinline{write_head()} 将头块写入磁盘：这是提交点，此写入之后的崩溃将导致恢复从日志中重放事务的写入。\indexcode{install_trans} 从日志中读取每个块并将其写入文件系统中的正确位置。最后，\lstinline{end_op} 写入一个计数为零的日志头；这必须在下一个事务开始写入日志块之前完成，这样崩溃就不会导致恢复将一个事务的头与后续事务的日志块一起使用。

\indexcode{recover_from_log} 从 \indexcode{initlog} 调用，而 \indexcode{initlog} 在引导期间，第一个用户进程运行之前，从 \indexcode{fsinit} 调用。它读取日志头，如果头指示日志包含一个已提交的事务，则模拟 \lstinline{end_op} 的操作。

日志的一个示例用法出现在 \indexcode{filewrite} 中。事务如下所示：
\begin{lstlisting}[]
      begin_op();
      ilock(f->ip);
      r = writei(f->ip, ...);
      iunlock(f->ip);
      end_op();
\end{lstlisting}
此代码被包装在一个循环中，该循环将大的写入分解为每次仅几个扇区的单个事务，以避免溢出日志。调用 \indexcode{writei} 会在此事务中写入许多块：文件的 inode、一个或多个位图块以及一些数据块。
%% 
%% 
%% 
\section{代码：块分配器}

文件和目录内容存储在磁盘块中，这些块必须从空闲池中分配。Xv6 的块分配器在磁盘上维护一个空闲位图，每个块对应一个位。零位表示对应的块是空闲的；一位表示它正在使用中。程序 \lstinline{mkfs} 设置与引导扇区、超级块、日志块、inode 块和位图块相对应的位。

块分配器提供两个函数：\indexcode{balloc} 分配一个新的磁盘块，\indexcode{bfree} 释放一个块。在 \lstinline{balloc} 中，从块 0 到 \lstinline{sb.size}（文件系统中的块数）循环考虑每个块。它寻找位图位为零的块，表示它是空闲的。如果 \lstinline{balloc} 找到这样的块，它更新位图并返回该块。为了提高效率，循环被分成两部分。外层循环读取每个位图块。内层循环检查单个位图块中的所有每块位数 (\lstinline{BPB})。如果两个进程同时尝试分配一个块可能发生的竞争，由于缓冲区缓存一次只允许一个进程使用任何一个位图块，因此被阻止。

\lstinline{bfree} 找到正确的位图块并清除正确的位。同样，\lstinline{bread} 和 \lstinline{brelse} 所隐含的独占使用避免了显式锁定的需要。

与本章其余部分描述的大部分代码一样，\lstinline{balloc} 和 \lstinline{bfree} 必须在事务内部调用。
%% 
%%  -------------------------------------------
%% 
\section{Inode 层}

术语 inode 可以有两个相关的含义。它可能指的是磁盘上的数据结构，包含文件的大小和数据块编号列表。或者，“inode”可能指的是内存中的 inode，它包含磁盘上 inode 的副本以及内核中需要的额外信息。

磁盘上的 inode 被压缩在磁盘上一个称为 inode 块的连续区域中。每个 inode 大小相同，因此给定一个数字 n，很容易找到磁盘上的第 n 个 inode。实际上，这个数字 n 被称为 inode 编号或 i-number，是实现在标识 inode 的方式。

磁盘上的 inode 由 \indexcode{struct dinode} 定义。\lstinline{type} 字段区分文件、目录和特殊文件（设备）。类型为零表示磁盘上的 inode 是空闲的。\lstinline{nlink} 字段统计引用此 inode 的目录项数量，以便识别何时应释放磁盘上的 inode 及其数据块。\lstinline{size} 字段记录文件内容的字节数。\lstinline{addrs} 数组记录保存文件内容的磁盘块的块号。

内核将一组活动 inode 保存在内存中一个称为 \indexcode{itable} 的表中；\indexcode{struct inode} 是磁盘上 \lstinline{struct dinode} 的内存副本。仅当有 C 指针引用该 inode 时，内核才将 inode 存储在内存中。\lstinline{ref} 字段统计引用内存中 inode 的 C 指针数量，如果引用计数降至零，内核将从内存中丢弃该 inode。\indexcode{iget} 和 \indexcode{iput} 函数获取和释放指向 inode 的指针，修改引用计数。指向 inode 的指针可以来自文件描述符、当前工作目录和瞬态内核代码，例如 \lstinline{kexec}。

在 xv6 的 inode 代码中有四种锁或类似锁的机制。\lstinline{itable.lock} 保护不变性：inode 在 inode 表中最多出现一次，以及内存中 inode 的 \lstinline{ref} 字段统计指向该 inode 的内存中指针数量。每个内存中的 inode 有一个包含睡眠锁的 \lstinline{lock} 字段，它确保对 inode 字段（如文件长度）以及 inode 的文件或目录内容块的独占访问。inode 的 \lstinline{ref} 如果大于零，会导致系统将该 inode 保留在表中，并且不将表项重用于不同的 inode。最后，每个 inode 包含一个 \lstinline{nlink} 字段（在磁盘上，如果在内存中则复制到内存中），它统计引用文件的目录项数量；如果链接计数大于零，xv6 不会释放 inode。

由 \lstinline{iget()} 返回的 \lstinline{struct inode} 指针保证在相应调用 \lstinline{iput()} 之前有效；inode 不会被删除，并且指针引用的内存不会重用于不同的 inode。\lstinline{iget()} 提供对 inode 的非独占访问，因此可以有许多指针指向同一个 inode。文件系统代码的许多部分依赖于 \lstinline{iget()} 的这种行为，既用于持有对 inode 的长期引用（如打开的文件和当前目录），也用于在操作多个 inode 的代码中（如路径名查找）防止竞争并避免死锁。

\lstinline{iget} 返回的 \lstinline{struct inode} 可能没有任何有用的内容。为了确保它包含磁盘上 inode 的副本，代码必须调用 \indexcode{ilock}。这会锁定 inode（因此没有其他进程可以 \lstinline{ilock} 它），并且如果尚未读取，则从磁盘读取 inode。\lstinline{iunlock} 释放 inode 上的锁。将获取 inode 指针与锁定分开有助于在某些情况下避免死锁，例如在目录查找期间。多个进程可以持有由 \lstinline{iget} 返回的指向 inode 的 C 指针，但一次只有一个进程可以锁定该 inode。

inode 表只存储内核代码或数据结构持有 C 指针的 inode。它的主要工作是同步多个进程的访问。inode 表也恰好缓存了频繁使用的 inode，但缓存是次要的；如果一个 inode 被频繁使用，缓冲区缓存很可能会将其保存在内存中。修改内存中 inode 的代码使用 \lstinline{iupdate} 将其写入磁盘。

%% 
%%  -------------------------------------------
%% 
\section{代码：Inodes}

为了分配一个新的 inode（例如，在创建文件时），xv6 调用 \indexcode{ialloc}。\lstinline{ialloc} 类似于 \indexcode{balloc}：它循环遍历磁盘上的 inode 结构，一次一个块，寻找标记为空闲的 inode。当找到一个时，它通过将新的 \lstinline{type} 写入磁盘来声明它，然后通过尾部调用 \indexcode{iget} 从 inode 表返回一个条目。\lstinline{ialloc} 的正确操作依赖于这样一个事实：一次只能有一个进程持有对 \lstinline{bp} 的引用：\lstinline{ialloc} 可以确信没有其他进程同时看到该 inode 可用并试图声明它。

\lstinline{iget} 在 inode 表中查找具有所需设备和 inode 编号的活动条目（\lstinline{ip->ref > 0}）。如果找到，它返回对该 inode 的新引用。当 \indexcode{iget} 扫描时，它记录第一个空槽的位置，如果它需要分配一个表项，则使用该空槽。

代码必须在读取或写入其元数据或内容之前使用 \indexcode{ilock} 锁定 inode。\lstinline{ilock} 为此目的使用睡眠锁。一旦 \indexcode{ilock} 独占访问 inode，如果需要，它会从磁盘（更可能从缓冲区缓存）读取 inode。函数 \indexcode{iunlock} 释放睡眠锁，这可能会唤醒任何正在休眠的进程。

\lstinline{iput} 通过递减引用计数来释放指向 inode 的 C 指针。如果这是最后一个引用，inode 表中该 inode 的槽位现在空闲，可以重用于不同的 inode。

如果 \indexcode{iput} 发现没有 C 指针引用一个 inode，并且该 inode 没有指向它的链接（不在任何目录中），那么该 inode 及其数据块必须被释放。\lstinline{iput} 调用 \indexcode{itrunc} 将文件截断为零字节，释放数据块；将 inode 类型设置为 0（未分配）；并将 inode 写入磁盘。

在释放 inode 的情况下，\indexcode{iput} 中的锁定协议值得仔细研究。一个危险是，并发的线程可能正在 \lstinline{ilock} 中等待使用这个 inode（例如，读取文件或列出目录），并且不会准备好发现 inode 不再被分配。这不可能发生，因为如果 inode 没有指向它的链接并且 \lstinline{ip->ref} 为 1，系统调用无法获得指向内存中 inode 的指针。那一个引用是调用 \lstinline{iput} 的线程拥有的引用。另一个主要危险是，对 \lstinline{ialloc} 的并发调用可能选择与 \lstinline{iput} 正在释放的相同的 inode。这只有在 \lstinline{iupdate} 写入磁盘使得 inode 类型为零之后才可能发生。这种竞争是良性的；分配线程会在读取或写入 inode 之前礼貌地等待获取 inode 的睡眠锁，而此时 \lstinline{iput} 已经完成了对该 inode 的操作。

\lstinline{iput()} 可以写入磁盘。这意味着任何使用文件系统的系统调用都可能写入磁盘，因为该系统调用可能是最后一个引用该文件的调用。即使是像 \lstinline{read()} 这样看似只读的调用，最终也可能调用 \lstinline{iput()}。这反过来意味着，即使是只读的系统调用，如果它们使用文件系统，也必须包装在事务中。

\lstinline{iput()} 和崩溃之间存在一个具有挑战性的交互。当文件的链接计数降至零时，\lstinline{iput()} 不会立即截断文件，因为某个进程可能仍然在内存中持有对该 inode 的引用：一个进程可能因为成功打开了文件而仍在读取和写入它。但是，如果在最后一个进程关闭该文件的文件描述符之前发生崩溃，那么该文件在磁盘上将被标记为已分配，但没有目录项指向它。

文件系统通过两种方式之一处理这种情况。简单的解决方案是，在恢复时，重启后，文件系统扫描整个文件系统，查找标记为已分配但没有目录项指向它们的文件。如果存在任何此类文件，则可以释放这些文件。

第二种解决方案不需要扫描文件系统。在这种解决方案中，文件系统在磁盘上（例如，在超级块中）记录链接计数降至零但引用计数不为零的文件的 inode 编号。如果文件系统在其引用计数达到 0 时删除该文件，则通过从列表中删除该 inode 来更新磁盘上的列表。恢复时，文件系统释放列表中的任何文件。

Xv6 两种解决方案都没有实现，这意味着 inode 可能在磁盘上被标记为已分配，即使它们不再使用。这意味着随着时间的推移，xv6 面临可能耗尽磁盘空间的风险。
%% 
%% 
%% 
\section{代码：Inode 内容}

\begin{figure}[t]
\center
\includegraphics[scale=0.5]{fig/inode.pdf}
\caption{文件在磁盘上的表示。}
\label{fig:inode}
\end{figure}

磁盘上的 inode 结构 \indexcode{struct dinode} 包含一个大小和一个块号数组（见图~\ref{fig:inode}）。inode 数据位于 \lstinline{dinode} 的 \lstinline{addrs} 数组中列出的块中。前 \indexcode{NDIRECT} 个数据块列在数组的前 \lstinline{NDIRECT} 个条目中；这些块称为直接块。接下来的 \indexcode{NINDIRECT} 个数据块没有列在 inode 中，而是列在一个称为间接块的数据块中。\lstinline{addrs} 数组的最后一个条目给出了间接块的地址。因此，文件的前 12 kB（\lstinline{NDIRECT x BSIZE}）字节可以从 inode 中列出的块加载，而接下来的 256 kB（\lstinline{NINDIRECT x BSIZE}）字节只能在查阅间接块后加载。这是一个很好的磁盘表示，但对客户端来说很复杂。函数 \indexcode{bmap} 管理表示，以便更高级别的例程（例如我们稍后将看到的 \indexcode{readi} 和 \indexcode{writei}）不需要管理这种复杂性。\lstinline{bmap} 返回 inode \lstinline{ip} 的第 \lstinline{bn} 个数据块的磁盘块号。如果 \lstinline{ip} 还没有这样的块，\lstinline{bmap} 会分配一个。

函数 \indexcode{bmap} 首先处理简单的情况：前 \indexcode{NDIRECT} 个块列在 inode 本身中。接下来的 \indexcode{NINDIRECT} 个块列在 \lstinline{ip->addrs[NDIRECT]} 的间接块中。\lstinline{bmap} 读取间接块，然后从块内的正确位置读取块号。如果块号超过 \lstinline{NDIRECT+NINDIRECT}，\lstinline{bmap} 会 panic；\lstinline{writei} 包含检查以防止这种情况发生。

\lstinline{bmap} 根据需要分配块。\lstinline{ip->addrs[]} 或间接条目为零表示没有分配块。当 \lstinline{bmap} 遇到零时，它会用按需分配的新块号替换它们。

\indexcode{itrunc} 释放文件的块，将 inode 的大小重置为零。\lstinline{itrunc} 首先释放直接块，然后释放间接块中列出的块，最后释放间接块本身。

\lstinline{bmap} 使 \indexcode{readi} 和 \indexcode{writei} 能够轻松访问 inode 的数据。\lstinline{readi} 首先确保偏移量和计数不超过文件末尾。从文件末尾之后开始的读取返回错误，而在文件末尾或跨越文件末尾的读取返回的字节数少于请求的字节数。主循环处理文件的每个块，将数据从缓冲区复制到 \lstinline{dst} 中。\indexcode{writei} 与 \indexcode{readi} 相同，但有三个例外：从文件末尾或跨越文件末尾的写入会增长文件，直到最大文件大小；循环将数据复制到缓冲区中而不是复制出来；如果写入扩展了文件，\indexcode{writei} 必须更新其大小。

函数 \indexcode{stati} 将 inode 元数据复制到 \lstinline{stat} 结构中，该结构通过 \indexcode{stat} 系统调用暴露给用户程序。
%% 
%% 
%% 
\section{代码：目录层}

目录在内部实现得与文件非常相似。其 inode 类型为 \indexcode{T_DIR}，其数据是一系列目录项。每个条目是一个 \indexcode{struct dirent}，包含一个名称和一个 inode 编号。名称最多为 \indexcode{DIRSIZ}（14）个字符；如果更短，则以 NULL（0）字节终止。inode 编号为零的目录项是空闲的。

函数 \indexcode{dirlookup} 在目录中搜索具有给定名称的条目。如果找到，它返回指向相应 inode 的指针（未锁定），并设置 \lstinline{*poff} 为目录中条目的字节偏移量，以防调用者希望编辑它。如果 \lstinline{dirlookup} 找到具有正确名称的条目，它会更新 \lstinline{*poff} 并返回通过 \indexcode{iget} 获得的未锁定 inode。\lstinline{dirlookup} 是 \lstinline{iget} 返回未锁定 inode 的原因。调用者已锁定 \lstinline{dp}，因此如果查找的是 \indexcode{.}（当前目录的别名），则在返回之前尝试锁定 inode 将尝试重新锁定 \lstinline{dp} 并导致死锁。（还有涉及多个进程和 \indexcode{..}（父目录的别名）的更复杂的死锁情况；\lstinline{.} 不是唯一的问题。）调用者可以解锁 \lstinline{dp}，然后锁定 \lstinline{ip}，确保一次只持有一个锁。

函数 \indexcode{dirlink} 将具有给定名称和 inode 编号的新目录项写入目录 \lstinline{dp} 中。如果名称已存在，\lstinline{dirlink} 返回错误。主循环读取目录项寻找未分配的条目。当找到时，它提前停止循环，并将 \lstinline{off} 设置为可用条目的偏移量。否则，循环结束，\lstinline{off} 设置为 \lstinline{dp->size}。无论哪种方式，\lstinline{dirlink} 然后通过在偏移量 \lstinline{off} 处写入来向目录添加一个新条目。
%% 
%% 
%% 
\section{代码：路径名}

路径名查找涉及一系列对 \indexcode{dirlookup} 的调用，每个路径分量一次。\lstinline{namei} 计算 \lstinline{path} 并返回相应的 \lstinline{inode}。函数 \indexcode{nameiparent} 是一个变体：它在最后一个元素之前停止，返回父目录的 inode 并将最后一个元素复制到 \lstinline{name} 中。两者都调用通用函数 \indexcode{namex} 来完成实际工作。

\lstinline{namex} 首先决定路径评估从何处开始。如果路径以斜杠开头，则从根目录开始评估；否则，从当前目录开始。然后它使用 \indexcode{skipelem} 依次考虑路径的每个元素。循环的每次迭代必须在当前 inode \lstinline{ip} 中查找 \lstinline{name}。迭代开始时锁定 \lstinline{ip} 并检查它是否是一个目录。如果不是，查找失败。（锁定 \lstinline{ip} 是必要的，不是因为它下面的 \lstinline{ip->type} 可以更改——它不能——而是因为在 \indexcode{ilock} 运行之前，不能保证 \lstinline{ip->type} 已从磁盘加载。）如果调用是 \indexcode{nameiparent} 并且这是最后一个路径元素，则根据 \lstinline{nameiparent} 的定义，循环提前停止；最后一个路径元素已被复制到 \lstinline{name} 中，因此 \indexcode{namex} 只需要返回未锁定的 \lstinline{ip}。最后，循环使用 \indexcode{dirlookup} 查找路径元素，并通过设置 \lstinline{ip = next} 为下一次迭代做准备。当循环用完路径元素时，它返回 \lstinline{ip}。

过程 \lstinline{namex} 可能需要很长时间才能完成：它可能涉及多个磁盘操作来读取路径名中遍历的目录的 inode 和目录块（如果它们不在缓冲区缓存中）。Xv6 被精心设计，以便如果一个内核线程对 \lstinline{namex} 的调用被磁盘 I/O 阻塞，另一个内核线程查找不同的路径名可以并发进行。\lstinline{namex} 单独锁定路径中的每个目录，以便在不同目录中的查找可以并行进行。

这种并发性带来了一些挑战。例如，当一个内核线程正在查找一个路径名时，另一个内核线程可能正在通过取消链接一个目录来更改目录树。一个潜在的风险是，查找可能正在搜索一个已被另一个内核线程删除的目录，并且其块已被重用于另一个目录或文件。

Xv6 避免了这种竞争。例如，当在 \lstinline{namex} 中执行 \lstinline{dirlookup} 时，查找线程持有目录的锁，并且 \lstinline{dirlookup} 返回一个使用 \lstinline{iget} 获得的 inode。\lstinline{iget} 增加 inode 的引用计数。只有在从 \lstinline{dirlookup} 接收到 inode 之后，\lstinline{namex} 才释放目录的锁。现在另一个线程可以从目录中取消链接该 inode，但 xv6 还不会删除该 inode，因为 inode 的引用计数仍然大于零。

另一个风险是死锁。例如，当查找 "." 时，\lstinline{next} 指向与 \lstinline{ip} 相同的 inode。在释放 \lstinline{ip} 上的锁之前锁定 \lstinline{next} 将导致死锁。为了避免这种死锁，\lstinline{namex} 在获取 \lstinline{next} 上的锁之前先解锁目录。这里我们再次看到将 \lstinline{iget} 和 \lstinline{ilock} 分开的重要性。
%% 
%% 
%% 
\section{文件描述符层}

Unix 接口的一个很酷的方面是，Unix 中的大多数资源都被表示为文件，包括控制台、管道等设备，当然还有真实的文件。文件描述符层是实现这种统一性的层。

如第~\ref{CH:UNIX}~章所述，Xv6 为每个进程提供了自己的打开文件表，即文件描述符。每个打开的文件由一个 \indexcode{struct file} 表示，它包装了一个 inode 或一个管道，以及一个 I/O 偏移量。每次调用 \indexcode{open} 都会创建一个新的打开文件（一个新的 \lstinline{struct file}）：如果多个进程独立打开同一个文件，不同的实例将有不同的 I/O 偏移量。另一方面，一个单独的打开文件（相同的 \lstinline{struct file}）可以在一个进程的文件表中多次出现，也可以在多个进程的文件表中出现。如果某个进程使用 \indexcode{open} 打开文件，然后使用 \indexcode{dup} 创建别名，或者使用 \indexcode{fork} 与子进程共享它，就会发生这种情况。引用计数跟踪对特定打开文件的引用数量。一个文件可以以只读、只写或读写方式打开。\lstinline{readable} 和 \lstinline{writable} 字段跟踪这一点。

系统中所有打开的文件都保存在一个全局文件表 \indexcode{ftable} 中。文件表具有分配文件（\indexcode{filealloc}）、创建重复引用（\indexcode{filedup}）、释放引用（\indexcode{fileclose}）以及读写数据（\indexcode{fileread} 和 \indexcode{filewrite}）的函数。

前三个遵循现在熟悉的形式。\lstinline{filealloc} 扫描文件表寻找一个未被引用的文件（\lstinline{f->ref == 0}）并返回一个新引用；\indexcode{filedup} 增加引用计数；\indexcode{fileclose} 减少引用计数。当文件的引用计数达到零时，\lstinline{fileclose} 根据类型释放底层的管道或 inode。

函数 \indexcode{filestat}、\indexcode{fileread} 和 \indexcode{filewrite} 实现对文件的 \indexcode{stat}、\indexcode{read} 和 \indexcode{write} 操作。\lstinline{filestat} 只允许用于 inode，并调用 \indexcode{stati}。\lstinline{fileread} 和 \lstinline{filewrite} 检查打开模式是否允许该操作，然后将调用传递给管道或 inode 的实现。如果文件代表一个 inode，\lstinline{fileread} 和 \lstinline{filewrite} 使用 I/O 偏移量作为操作的偏移量，然后前进它。管道没有偏移量的概念。回想一下，inode 函数要求调用者处理锁定。inode 锁定有一个方便的副作用，即读取和写入偏移量是原子更新的，因此同时写入同一文件的多个写入不会覆盖彼此的数据，尽管它们的写入可能会交错。
%% 
%% 
%% 
\section{代码：系统调用}

有了较低层提供的函数，大多数系统调用的实现都很简单（参见 \fileref{kernel/sysfile.c}）。有几个调用值得仔细研究。

函数 \indexcode{sys_link} 和 \indexcode{sys_unlink} 编辑目录，创建或删除对 inode 的引用。它们很好地展示了使用事务的能力。\lstinline{sys_link} 首先获取其两个字符串参数 \lstinline{old} 和 \lstinline{new}。假设 \lstinline{old} 存在且不是一个目录，\lstinline{sys_link} 增加其 \lstinline{ip->nlink} 计数。然后 \lstinline{sys_link} 调用 \indexcode{nameiparent} 查找 \lstinline{new} 的父目录和最后一个路径元素，并创建一个指向 \lstinline{old} 的 inode 的新目录项。新的父目录必须存在，并且与现有 inode 位于同一设备上：inode 编号仅在单个磁盘上具有唯一含义。如果发生这样的错误，\indexcode{sys_link} 必须返回并减少 \lstinline{ip->nlink}。

事务简化了实现，因为它需要更新多个磁盘块，但我们不必担心执行它们的顺序。它们要么全部成功，要么全部不成功。例如，如果没有事务，在创建链接之前更新 \lstinline{ip->nlink} 将使文件系统暂时处于不安全状态，并且在两者之间发生崩溃可能会造成严重破坏。有了事务，我们就不必担心这个问题。

\lstinline{sys_link} 为现有的 inode 创建一个新名称。函数 \indexcode{create} 为一个新的 inode 创建一个新名称。它是三个文件创建系统调用的泛化：带有 \indexcode{O_CREATE} 标志的 \indexcode{open} 创建一个新的普通文件，\indexcode{mkdir} 创建一个新目录，\indexcode{mkdev} 创建一个新设备文件。与 \indexcode{sys_link} 一样，\indexcode{create} 首先调用 \indexcode{nameiparent} 获取父目录的 inode。然后它调用 \indexcode{dirlookup} 检查名称是否已经存在。如果名称存在，\lstinline{create} 的行为取决于它用于哪个系统调用：\lstinline{open} 的语义与 \indexcode{mkdir} 和 \indexcode{mkdev} 不同。如果 \lstinline{create} 代表 \lstinline{open}（\lstinline{type == T_FILE}）使用，并且存在的名称本身是一个普通文件，那么 \lstinline{open} 将其视为成功，因此 \lstinline{create} 也如此。否则，这是一个错误。如果名称不存在，\lstinline{create} 现在使用 \indexcode{ialloc} 分配一个新的 inode。如果新 inode 是一个目录，\lstinline{create} 用 \indexcode{.} 和 \indexcode{..} 条目初始化它。最后，在数据正确初始化之后，\indexcode{create} 可以将其链接到父目录中。与 \indexcode{sys_link} 一样，\lstinline{create} 同时持有两个 inode 锁：\lstinline{ip} 和 \lstinline{dp}。没有死锁的可能性，因为 inode \lstinline{ip} 是新分配的：系统中的任何其他进程都不会持有 \lstinline{ip} 的锁并尝试锁定 \lstinline{dp}。

使用 \lstinline{create}，很容易实现 \indexcode{sys_open}、\indexcode{sys_mkdir} 和 \indexcode{sys_mknod}。\lstinline{sys_open} 是最复杂的，因为创建新文件只是它所能做的一小部分。如果传递给 \indexcode{open} 的是 \indexcode{O_CREATE} 标志，它调用 \lstinline{create}。否则，它调用 \indexcode{namei}。\lstinline{create} 返回一个锁定的 inode，但 \lstinline{namei} 不会，因此 \indexcode{sys_open} 必须自己锁定 inode。这为检查目录是否仅以只读方式打开提供了一个方便的地方。假设以某种方式获得了 inode，\lstinline{sys_open} 分配一个文件和一个文件描述符，然后填充文件。请注意，没有其他进程可以访问部分初始化的文件，因为它只在当前进程的表中。

第~\ref{CH:SLEEP}~章在我们还没有文件系统之前就检查了管道的实现。函数 \indexcode{sys_pipe} 通过提供创建管道对的方法将该实现连接到文件系统。其参数是一个指向两个整数空间的指针，它将在其中记录两个新文件描述符。然后它分配管道并安装文件描述符。
%% 
%%  -------------------------------------------
%% 
\section{现实世界}

现实世界操作系统中的缓冲区缓存比 xv6 的复杂得多，但它服务于相同的两个目的：缓存和同步对磁盘的访问。Xv6 的缓冲区缓存（如 V6 的）使用简单的最近最少使用（LRU）驱逐策略；可以实施许多更复杂的策略，每个策略对某些工作负载有利，而对其他工作负载不那么有利。更高效的 LRU 缓存将消除链表，而是使用哈希表进行查找，使用堆进行 LRU 驱逐。现代缓冲区缓存通常与虚拟内存系统集成以支持内存映射文件。

Xv6 的日志系统效率低下。提交不能与文件系统系统调用同时发生。系统记录整个块，即使一个块中只有几个字节被更改。它一次一个块地执行同步日志写入，每个写入可能需要整个磁盘旋转时间。真实的日志系统解决了所有这些问题。

日志记录不是提供崩溃恢复的唯一方法。早期的文件系统在重启期间使用清理程序（例如 UNIX \indexcode{fsck} 程序）检查每个文件和目录以及块和 inode 空闲列表，查找并解决不一致性。对于大型文件系统，清理可能需要数小时，并且在某些情况下，无法以使得原始系统调用具有原子性的方式解决不一致性。从日志中恢复要快得多，并且在崩溃时使系统调用具有原子性。

Xv6 使用与早期 UNIX 相同的基本磁盘 inode 和目录布局；这种方案多年来一直非常持久。BSD 的 UFS/FFS 和 Linux 的 ext2/ext3 使用基本相同的数据结构。文件系统布局中最低效的部分是目录，它在每次查找时需要线性扫描所有磁盘块。当目录只有几个磁盘块时，这是合理的，但对于包含大量文件的目录来说，这是昂贵的。微软 Windows 的 NTFS、macOS 的 HFS 和 Solaris 的 ZFS 等，将目录实现为磁盘上的平衡树块。这很复杂，但保证了对数时间的目录查找。

Xv6 对磁盘故障很天真：如果磁盘操作失败，xv6 会 panic。这是否合理取决于硬件：如果操作系统位于使用冗余来掩盖磁盘故障的特殊硬件之上，也许操作系统看到故障的频率如此之低，以至于 panic 是可以的。另一方面，使用普通磁盘的操作系统应该预期故障并更优雅地处理它们，这样一个文件中一个块的丢失不会影响文件系统其余部分的使用。

Xv6 要求文件系统适合一个磁盘设备并且大小不变。随着大型数据库和多媒体文件推动存储需求不断增长，操作系统正在开发消除“每个文件系统一个磁盘”瓶颈的方法。基本方法是将许多磁盘组合成一个逻辑磁盘。硬件解决方案如 RAID 仍然是最流行的，但当前的趋势是尽可能多地将这种逻辑在软件中实现。这些软件实现通常允许丰富的功能，例如通过动态添加或移除磁盘来增长或缩小逻辑设备。当然，一个可以动态增长或缩小的存储层需要一个可以做到这一点的文件系统：xv6 使用的固定大小的 inode 块数组在这样的环境中无法很好地工作。将磁盘管理与文件系统分开可能是最清晰的设计，但两者之间的复杂接口导致一些系统（如 Sun 的 ZFS）将它们结合起来。

Xv6 的文件系统缺少现代文件系统的许多其他功能；例如，它不支持快照和增量备份。

现代 Unix 系统允许使用与磁盘存储相同的系统调用访问许多种类的资源：命名管道、网络连接、远程访问的网络文件系统，以及诸如 \lstinline{/proc} 之类的监视和控制接口。这些系统通常不采用 xv6 的 \indexcode{fileread} 和 \indexcode{filewrite} 中的 \lstinline{if} 语句，而是为每个打开的文件提供一个每个操作一个函数指针的表，并调用函数指针来调用该 inode 对该调用的实现。网络文件系统和用户级文件系统提供将这些调用转换为网络 RPC 并等待响应后返回的函数。
%% 
%%  -------------------------------------------
%% 
\section{练习}

\begin{enumerate}

\item 为什么在 \lstinline{balloc} 中 panic？ xv6 能恢复吗？

\item 为什么在 \lstinline{ialloc} 中 panic？ xv6 能恢复吗？

\item 为什么 \lstinline{filealloc} 在文件用完时不会 panic？为什么这更常见因此值得处理？

\item 假设对应于 \lstinline{ip} 的文件在 \lstinline{sys_link} 调用 \lstinline{iunlock(ip)} 和 \lstinline{dirlink} 之间被另一个进程取消链接。链接会被正确创建吗？为什么或为什么不？

\item \lstinline{create} 进行了四个函数调用（一次 \lstinline{ialloc} 和三次 \lstinline{dirlink}），它要求这些调用成功。如果有任何一个失败，\lstinline{create} 调用 \lstinline{panic}。为什么这是可以接受的？为什么这四个调用中的任何一个都不会失败？

\item \lstinline{sys_chdir} 在 \lstinline{iput(cp->cwd)} 之前调用 \lstinline{iunlock(ip)}，后者可能尝试锁定 \lstinline{cp->cwd}，但将 \lstinline{iunlock(ip)} 推迟到 \lstinline{iput} 之后不会导致死锁。为什么？

\item 实现 \lstinline{lseek} 系统调用。支持 \lstinline{lseek} 还需要你修改 \lstinline{filewrite}，以便如果 \lstinline{lseek} 将 \lstinline{off} 设置到 \lstinline{f->ip->size} 之后，用零填充文件中的空洞。

\item 向 \lstinline{open} 添加 \lstinline{O_TRUNC} 和 \lstinline{O_APPEND}，以便 shell 中的 \lstinline{>} 和 \lstinline{>>} 运算符能够工作。

\item 修改文件系统以支持符号链接。

\item 修改文件系统以支持命名管道。

\item 修改文件和 VM 系统以支持内存映射文件。

\end{enumerate}
